% AI Context (for automated fix generation)
\needspace{8cm}
\section{AI Context}

% This section provides structured technical details that AI coding
% assistants can consume directly to produce working fixes.
% Only sections with populated data should be included in the final report.

\subsection{Context Summary}

% High-level summary of the technical context surrounding this defect.

\begin{keyinsight}
Provide a concise summary of the technical context: what subsystem is
affected, what the failure mode is, and what approach is recommended
for resolution.
\end{keyinsight}

\subsection{Affected Files}

% List the source files relevant to this defect.

\begin{itemize}
    \item \texttt{path/to/affected/file.ext}
    \item \texttt{path/to/another/file.ext}
\end{itemize}

\subsection{Code Snippets}

% Include the relevant code that exhibits or contributes to the defect.

\begin{wyrdcode}[Relevant Code]{text}
// Paste the relevant code snippet here
\end{wyrdcode}

\subsection{Stack Trace}

% Include the stack trace captured when the defect occurred.

\begin{wyrdcode}[Stack Trace]{text}
// Paste the stack trace here
\end{wyrdcode}

\subsection{Suggested Fix}

% Describe or show the recommended code change to resolve the defect.

\begin{technicaldetail}
Describe the suggested fix in precise technical terms, referencing
the specific code changes needed.
\end{technicaldetail}

\subsection{Related Commits}

% Reference any commits that introduced or are related to this defect.

\begin{wyrdtab}{l l}
\textbf{Commit SHA} & \textbf{Description} \\
abc1234             & Brief commit description \\
\end{wyrdtab}

\subsection{Fix Confidence}

% Indicate the confidence level of the suggested fix.
% Score: 0--100 where higher values indicate greater confidence.

\begin{callout}
Fix confidence score: \textbf{N/100}. A higher score indicates
the suggested fix is more likely to resolve the defect without
introducing regressions.
\end{callout}
